\documentclass[12pt, xelatex,ja=standard,jafont=noto]{bxjsreport}

\usepackage{mxd_bachelorthesis}
\usepackage{natbib}
\usepackage{graphicx}
\usepackage{amsmath}
\usepackage{hyperref}
\hypersetup{
    colorlinks=false
}
\CJKspace

%----
\year{2019} %卒業年度
\date{2020年3月} %提出年月
\affiliation{立命館大学}
\department{情報理工学部情報理工学科}
\advisor{野間 春生 教授}
\title{非ユークリッド空間における\\球の充填と被覆の問題について}
\author{松村 耕平}

%----本文はここから---

\begin{document}
\maketitle

\begin{abstract}
本論文では ...
本論文では ...
本論文では ...
本論文では ...本論文では ...本論文では ...本論文では ...本論文では ...本論文では ...本論文では ...本論文では ...本論文では ...本論文では ...本論文では ...本論文では ...本論文では ...本論文では ...本論文では ...本論文では ...本論文では ...本論文では ...本論文では ...本論文では ...本論文では ...本論文では ...本論文では ...本論文では ...本論文では ...本論文では ...本論文では ...本論文では ...本論文では ...本論文では ...本論文では ...本論文では ...本論文では ...本論文では ...本論文では ...本論文では ...本論文では ...本論文では ...本論文では ...本論文では ...本論文では ...本論文では ...本論文では ...本論文では ...本論文では ...本論文では ...本論文では ...本論文では ...本論文では ...本論文では ...本論文では ...本論文では ...本論文では ...本論文では ...本論文では ...本論文では ...本論文では ...本論文では ...本論文では ...本論文では ...本論文では ...本論文では ...本論文では ...本論文では ...本論文では ...本論文では ...本論文では ...本論文では ...本論文では ...本論文では ...本論文では ...本論文では ...本論文では ...本論文では ...本論文では ...本論文では ...本論文では ...本論文では ...本論文では ...本論文では ...本論文では ...本論文では ...本論文では ...本論文では ...
\end{abstract}

%\begin{eabstract}
%This paper describes ...
%\end{eabstract}

\tableofcontents


\chapter{はじめに}
ぞうの卵はおいしいぞう。ぞうの卵はおいしいぞう。ぞうの卵はおいしいぞう。ぞうの卵はおいしいぞう。
\section{セクション1}
ぞうの卵はおいしいぞう。ぞうの卵はおいしいぞう。ぞうの卵はおいしいぞう。ぞうの卵はおいしいぞう。
\subsection{サブセクション1}
ぞうの卵はおいしいぞう。ぞうの卵はおいしいぞう。ぞうの卵はおいしいぞう。ぞうの卵はおいしいぞう。
\subsubsection{サブサブセクション1}
ぞうの卵はおいしいぞう。ぞうの卵はおいしいぞう。ぞうの卵はおいしいぞう。ぞうの卵はおいしいぞう。
\paragraph{パラグラフ}
図\ref{fig:lena}に示したように、ぞうの卵はおいしいぞう。ぞうの卵はおいしいぞう。ぞうの卵はおいしいぞう。ぞうの卵はおいしいぞう。

\begin{figure}[htbp]
\centering
\includegraphics[scale=0.2]{images/lena.jpg}
\caption{A standard test image widely used in the field of image processing.}
\label{fig:lena}
\end{figure}

\chapter{関連研究}
ぞうの卵はおいしいぞう。ぞうの卵はおいしいぞう。ぞうの卵はおいしいぞう。ぞうの卵はおいしいぞう\citep{6926269,Matsumura:2013:CAR:2468356.2468600, Matsumura:2014:CFS:2582051.2582068,Matsumura:2014:YTW:2667317.2667417,matsumura2011role,matsumura2012office}。
\chapter{アプローチ}
ぞうの卵はおいしいぞう。ぞうの卵はおいしいぞう。ぞうの卵はおいしいぞう。ぞうの卵はおいしいぞう。
\chapter{実験}
ぞうの卵はおいしいぞう。ぞうの卵はおいしいぞう。ぞうの卵はおいしいぞう。ぞうの卵はおいしいぞう。
\chapter{議論}
ぞうの卵はおいしいぞう。ぞうの卵はおいしいぞう。ぞうの卵はおいしいぞう。ぞうの卵はおいしいぞう。
\chapter{結論}
ぞうの卵はおいしいぞう。ぞうの卵はおいしいぞう。ぞうの卵はおいしいぞう。ぞうの卵はおいしいぞう。


% 参考文献はbibtexで
\bibliographystyle{plain}
\bibliography{references}

\chapter*{謝辞}
お父さん、お母さん、今まで育ててくれてありがとう。ぞうの卵はおいしいぞう。ぞうの卵はおいしいぞう。ぞうの卵はおいしいぞう。ぞうの卵はおいしいぞう。
\end{document}
